\subsection{Anillos excelentes y resolución de singularidades.}
\label{subsection:IV.2.7.9-fr}

\begin{env}[7.9.1]
\label{IV.2.s7.9.1-fr}
Dado un preesquema \emph{localmente noetheriano reducido} $X$, se llama
\emph{morfismo de desingularización} de $X$ a un morfismo $f:X'\to X$
tal que $X'$ sea \emph{regular} y $f$ sea \emph{propio} y \emph{birracional}.
Cuando existe tal morfismo, se dice que se pueden \emph{resolver las
singularidades de} $X$ (o, más simplemente, que se puede \emph{resolver}
$X$). Por ejemplo, si $A$ es un \emph{anillo japonés de dimensión} $1$,
se puede resolver $X=\operatorname{Spec}(A)$, pues el morfismo $X'\to X$
(donde $X'$ es la \emph{normalización} de $X$) es finito (y, por tanto,
propio) y birracional, y los anillos locales de $X'$ son anillos
íntegramente cerrados de dimensión a lo sumo $1$, por lo que son cuerpos o anillos de
valoración discreta y, en consecuencia, son regulares.
\end{env}
